\GuideSpread{What Strong Artin asks us to construct}{A global correspondence, with local parameters retained}{01_interface.pdf}{
Start with an irreducible finite image complex representation $\rho:G_F\to\GL_n(\C)$. Strong Artin asks for a cuspidal automorphic representation with the corresponding local parameters. Degree one is supplied by class field theory. The trivial character's $L$ function has its familiar pole.\par
The unramified polynomial is one visible interface. Matching it almost everywhere is weak attachment. All remaining excluded finite places and the archimedean places still require their own comparisons for full attachment.\par
The dashed arrow is the requested construction. This figure does not supply it. The rest of the guide follows exact reductions and local algebra around that arrow.
}{Preserve the global object and the normalization of every local test.}

\GuideSpread{A solvable tower has two different steps}{Clifford theory separates induction from invariant extension}{02_clifford.pdf}{
At a prime cyclic quotient $G/H\cong C_p$, an irreducible constituent of the restriction either moves in an orbit of size $p$ or is invariant.\par
The moving case is an induced representation. On an accessible prime cyclic tower, the established automorphic induction machinery mirrors that step under its hypotheses.\par
The invariant case produces an extension problem. Cyclic base change may supply descents of the automorphic restriction, but one must still attach the selected descent to the selected Galois extension.\par
The diagram concerns this solvable recursion. It is not a reduction of every finite group to cyclic quotients.
}{Induction and attached descent carry different mathematical burdens.}

\GuideSpread{Two torsors share a rotation law}{Labels are relative to a chosen origin}{03_torsor.pdf}{
For an invariant irreducible restriction, the $p$ complex extensions are $\rho_0\otimes\xi^j$. Under cyclic descent hypotheses, the descents are $\pi_0\otimes\epsilon^j$.\par
Both rows are acted on freely by $\Z/p\Z$. Choosing an origin labels each row. It does not prove that the two origins are attached.\par
The figure uses $p=3$ and a hypothetical offset 1. The gold correspondence is a possible alignment, not an automorphic construction. A simultaneous rotation changes the labels and preserves the same relative offset.
}{A torsor has a rotation action before it has a preferred origin.}

\GuideSpread{An intertwiner must close the cyclic relation}{Normalize its scalar $p$th power}{04_intertwiner.pdf}{
An invariant irreducible type gives an intertwiner $S$. After accounting for the chosen lift of the quotient generator, Schur's lemma leaves a scalar discrepancy in the cyclic relation.\par
Over $\C$, that nonzero scalar has a $p$th root. Rescaling $S$ removes the discrepancy. The remaining $p$ choices are the quotient character twists.\par
The exact matrix model has $p=3$, $S=2\operatorname{diag}(1,\zeta,\zeta^2)$, and $S^3=8I$. Dividing by 2 gives $B^3=I$. Multiplying $B$ by $\zeta^j$ gives three normalized choices.\par
This linear algebra illustrates extension normalization. Arithmetic attachment is a separate step.
}{Normalization supplies extensions and preserves their finite ambiguity.}

\GuideSpread{Finite Satake data form a finite alphabet}{A root bound produces a polynomial bound}{05_satake.pdf}{
At the unramified comparison places outside one fixed exceptional set, if the lower Artin roots lie in $\mu_M$ and cyclic base change has degree $d$, the local root argument puts normalized descent roots in $\mu_{dM}$.\par
A degree $n$ polynomial is determined by a multiset of $n$ roots. Thus at most $\binom{dM+n-1}{n}$ polynomials can occur. This is a bound on possible local data, not a count of automorphic representations.\par
The worked alphabet takes $M=4$, $d=3$, and $n=2$. There are 12 root choices and 78 unordered pairs with repetition. The circle shows all 12 roots and the triangular array counts all 78 multisets.
}{Finitely many local polynomials do not by themselves construct a global representation.}

\GuideSpread{Every coefficient has its own rotation degree}{Exterior powers expose the scalar action}{06_coefficients.pdf}{
Write $\det(TI-A)=\sum_{k=0}^n(-1)^k e_k(A)T^{n-k}$. Scalar multiplication gives $e_k(\zeta^j A)=\zeta^{jk}e_k(A)$.\par
The figure displays the exponents modulo 3 for degrees 1, 2, and 3. The trace rotates once. The second coefficient rotates twice. The determinant does not rotate at rank 3.\par
A coefficient can detect a label only when it is nonzero and its degree is not divisible by $p$. The sign convention in the polynomial matters when exact coefficients are compared.
}{An active coefficient carries both its value and its degree.}

\GuideSpread{Polynomial stabilizers come from active degrees}{The example $T^3-1$ is completely rotation blind}{07_stabilizer.pdf}{
For prime $p$, the stabilizer of the characteristic polynomial is the set of $\zeta\in\mu_p$ with $\zeta^k=1$ at every active degree $k$.\par
For $A=\operatorname{diag}(1,\zeta,\zeta^2)$, the polynomial is $T^3-1$. Only degree 3 is active. Multiplying the roots by $\zeta$ merely permutes the root multiset.\par
Contrast $\operatorname{diag}(1,1,\zeta)$. Its trace is $2+\zeta\ne0$, so its three scalar twists have different polynomials.\par
This coefficient stabilizer concerns a matrix polynomial. Extension spectrum stabilizers concern multiplicity labels and use a different object.
}{A vanished coefficient supplies no separator, even at a favorable degree.}

\GuideSpread{A determinant can miss a visible label}{An exact rank 3 example at the quotient prime 3}{08_blindness.pdf}{
Take $A=\operatorname{diag}(1,1,\zeta)$ with $\zeta^3=1$ and $\zeta\ne1$. Its elementary coefficients are $2+\zeta$, $1+2\zeta$, and $\zeta$.\par
The trace and second coefficient rotate with the quotient label. The determinant stays $\zeta$ because multiplying a rank 3 matrix by $\zeta$ multiplies its determinant by $\zeta^3=1$.\par
This is the difficult rank divisibility pattern. When $p\nmid n$, determinant data selects a candidate in an already known twist torsor. Candidate selection still requires attachment evidence for the other local parameters.
}{Determinant blindness is precise and does not imply blindness of all coefficients.}

\GuideSpread{Extension spectra retain multiplicity labels}{A stable type records more than its total multiplicity}{09_spectrum.pdf}{
For a stable irreducible $H$ type $\tau$, choose one extension. Its multiplicity space decomposes into quotient characters. Record it as $P_{\tau,V}(z)=\sum_a m_{\tau,a}z^a$ in $\Z[z]/(z^p-1)$.\par
Twisting by $\xi^j$ multiplies this polynomial by $z^j$. The displayed multiplicities $(2,1,0)$ rotate cyclically. Their sum stays 3, while the label distribution changes.\par
A moving orbit contributes an induced object and is unchanged by quotient twisting. The flat row $(1,1,1)$ illustrates rotational invariance of a regular quotient character packet, not an identification of stable and moving types.
}{Stable extension spectra preserve the labels that total multiplicity forgets.}

\GuideSpread{A global label must pass every local test}{Exact alignment and coarse filters have different geometry}{10_alignment.pdf}{
The global criterion is $\bigcap_v J_v\ne\varnothing$, where $J_v$ is the set of quotient shifts giving exact local agreement. In a prime cyclic quotient, an exact $J_v$ is empty, a singleton, or the whole quotient.\par
The upper example has $J_v=J_w=\{1\}$ and selects shift 1. Two exact singleton tests $\{1\}$ and $\{2\}$ would instead be incompatible.\par
The lower example uses coarse filters $\{0,1\}$, $\{1,2\}$, and $\{0,2\}$. Every pair intersects, but their total intersection is empty. These are arbitrary candidate filters, not exact prime cyclic alignment sets.
}{Check one common offset. Do not replace exact local sets by looser candidate filters.}

\GuideSpread{The total induced object forgets the origin}{Rotating a full packet only permutes its constituents}{11_induction.pdf}{
For a quotient character group $X$, $\bigoplus_{\chi\in X}V\otimes\chi\cong\Ind_H^G\operatorname{Res}_H V$. The automorphic induction identity has the same permutation symmetry in its established domain.\par
The finite model takes a one dimensional extension with label $j$. Summing all three twists gives one copy of each cube root for every $j$. The rows have different origins and identical total multiplicities.\par
Agreement of the total packet therefore cannot tell us which individual constituent matches the target. A product of local factors has the same loss of origin.
}{A norm identity does not select one factor of the norm.}

\GuideSpread{The quotient prime collapses all labels}{Congruence at 3 and separation at 7}{12_collapse.pdf}{
At a coefficient place above $p$, every $p$th root of unity reduces to 1. The quotient twists consequently have the same semisimple residual data.\par
For $p=3$, the three roots collapse in residue characteristic 3. The polynomial $(T-1)^3$ becomes $T^3-1$. In $\FF_7$, the cube roots $1,2,4$ remain distinct.\par
The table reduces the exact three rank 3 polynomials from the previous example. Their coefficients all agree at 3 and separate at 7.\par
Residual agreement does not prove ordinary equality. It also does not construct a residual Galois carrier from a Hecke eigensystem.
}{A collapsing prime gives congruence, while a separating prime retains the label.}

\GuideSpread{A finite alphabet has finitely many bad primes}{Distinct coefficient vectors survive almost every reduction}{13_separation.pdf}{
For each pair of distinct integral polynomials, choose a nonzero coefficient difference. Only finitely many coefficient places divide one of those differences. All other places separate the alphabet.\par
The integer example uses $(T-1)^2$ and $(T+1)^2$. Their middle coefficients differ by 4. They coincide modulo 2 and differ modulo 3, 5, and 7.\par
A good prime for the cyclic attachment problem must also avoid the quotient prime and preserve the absolute irreducibility of the restricted Galois representation.\par
The existence of good primes is established. A Galois realization at a good prime remains an additional input.
}{Reduction can recover equality only after injectivity on the entire relevant alphabet is proved.}

\GuideSpread{Trace words can encode a finite matrix group}{The bound needs semisimplicity and characteristic above the rank}{14_trace_bound.pdf}{
Let $\Gamma\subset\GL_n(k)$ be finite, with $k$ algebraically closed, characteristic $\ell>n$, and semisimple defining module. Its image algebra has a nondegenerate trace pairing.\par
Choose $d\le n^2$ group matrices spanning that algebra. The word of traces $\bigl(\Tr(\gamma\gamma_i)\bigr)_i$ determines $\gamma$. If there are $t$ trace values, $|\Gamma|\le t^d\le t^{n^2}$.\par
The exact model is the order 3 diagonal group over $\FF_7$, viewed over its algebraic closure. It has two trace values and $d=2$. Three distinct trace words fit inside four available words. The universal rank bound gives 16.
}{The trace pairing turns a finite alphabet into an injective encoding under explicit hypotheses.}

\GuideSpread{Pointwise spectra do not enforce multiplication}{A finite polynomial assignment that is not a representation}{15_multiplication.pdf}{
This is a finite algebra model, not an automorphic example. Assign $I$ to the identity of $C_2$ and $A=\operatorname{diag}(\zeta,\zeta^{-1})$ to its generator.\par
Both matrices have finite order, and their eigenvalues lie in $\mu_3$. Thus the assigned local polynomial alphabet is finite.\par
But the generator of $C_2$ must square to the identity, whereas $A^2\ne I$. No group representation has this assignment. The bottom row displays the exact failure over $\FF_7$ with $\zeta=2$.\par
This control isolates why local values need a globally compatible representation or pseudocharacter.
}{An alphabet constrains values. A representation also satisfies every multiplication relation.}

\GuideSpread{Finite traces force finite semisimple monodromy}{One l adic realization is an input, not an output of the alphabet}{16_monodromy.pdf}{
For a continuous semisimple Galois representation over a finite extension of $\Q_\ell$, unramified outside a finite set, finitely many Frobenius traces on a density one set force finitely many image traces. In the reductive Zariski closure, the connected component then has constant trace.\par
Its torus weight characters cannot cancel a nonconstant positive multiplicity character sum in characteristic zero. Every weight is trivial, so the faithful connected monodromy is trivial. The image is finite.\par
Semisimplicity matters. The lower model is the continuous unipotent representation of $\Z_\ell$ given by $t\mapsto\begin{pmatrix}1&t\\0&1\end{pmatrix}$. Its trace is always 2 and its image is infinite.\par
The theorem begins with a realization already constructed.
}{A rigidity theorem upgrades an existing semisimple realization.}

\GuideSpread{Cyclic $H^1$ keeps the norm condition}{A full matrix cohomology group need not be a line}{17_h1.pdf}{
For a cyclic group of order $q$ acting by $T$ on $M$, the standard formula is $H^1=\ker N/\im(T-I)$, with $N=I+T+\cdots+T^{q-1}$.\par
Over $\FF_3$, let $T$ be conjugation by $\operatorname{diag}(1,-1)$ on $2\times2$ matrices. Diagonal coordinates are fixed. Off diagonal coordinates change sign.\par
For $C_6$, the norm is zero. The quotient retains both diagonal coordinates, giving dimension 2. For $C_2$, the norm kills exactly the off diagonal coordinates, which are also the image of $T-I$, giving dimension 0.\par
Neither calculation identifies a selected continuous local Tate line.
}{Name the acting group, coefficient module, norm, and chosen local subquotient.}

\GuideSpread{Square zero defects need an image test}{Compatibility of an error is weaker than correctability}{18_h2.pdf}{
Across a square zero kernel, the multiplicativity error of a chosen representation lift satisfies a cocycle compatibility relation. Changing the lift changes the error by a correction map.\par
A correction exists exactly when the error class vanishes in the appropriate obstruction quotient. The left finite model has $J(a)=(a,0)$ and $B=0$ over $\FF_3$. Every error is compatible, but a vertical error is not in $\im J$.\par
The right example lifts $\begin{pmatrix}1&1\\0&1\end{pmatrix}$ from $\FF_2$ to $\Z/4$. Its square has an off diagonal error 2. Adding $2\operatorname{diag}(1,0)$ corrects it.\par
These models do not construct the required arithmetic $H^2$ comparison.
}{A cocycle relation locates an obstruction. Vanishing of its class supplies correction.}

\GuideSpread{A moving syzygy follows the changing point}{Correct the original system rather than an enlarged one}{19_syzygy.pdf}{
For $F(x)=(x,x^2)$, the exact relation is $(-x,1)F(x)=0$. At $x=2$ modulo 8, the error is $(2,4)$. The frozen residual row $(0,1)$ does not annihilate that whole error.\par
At the leading graded layer, the error divided by 2 reduces to $(1,0)$ modulo 2 and is in the Jacobian image. This is the layer on which the moving syzygy argument works.\par
The lower control takes $x^2+1$ modulo 4. It has no root. Adding an auxiliary coordinate $y$ gives a root $(1,2)$ of $x^2+1+y=0$, but $y\ne0$. The enlarged point does not solve the original equation.
}{Prove moving identities and residual exactness for the actual target equations.}

\GuideSpread{A scalar Newton correction gains valuation}{Exact rational iterates for $x^2=6$ at 5}{20_hensel_scalar.pdf}{
Begin at $x_0=1$ and use $x_{r+1}=x_r-(x_r^2-6)/(2x_r)$. The displayed values are exact rational numbers with 5 adic unit denominators.\par
The residual valuations are 1, 2, 4, 8, and 16. The identity $x_{r+1}^2-6=(x_r^2-6)^2/(4x_r^2)$ explains the doubling, since every $x_r$ remains a 5 adic unit.\par
The plotted five iterates are bounded computations. The identity and completeness give the usual local convergence under these hypotheses.\par
This scalar example illustrates valuation gain. A matrix power requires its ordered derivative, not a scalar substitute.
}{The error improves because the exact derivative and remainder satisfy the required estimates.}

\GuideSpread{A finite inverse certificate decides a matrix ball}{The square map at a unipotent matrix over $\Q_2$}{21_matrix_hensel.pdf}{
Let $A=\begin{pmatrix}1&1\\0&1\end{pmatrix}$. The square derivative is $D_A(X)=AX+XA$. With row major coordinates, the displayed integral matrix $B$ satisfies $BD=DB=4I$. Thus $Q=B/4$ is the exact inverse.\par
Here $s=2$ and choose $m=3>s$. The local criterion is that $B\operatorname{vec}(U-A^2)$ vanish modulo $2^{m+s}=32$.\par
The target $U=\begin{pmatrix}9&2\\16&9\end{pmatrix}$ passes and is the square of $F=\begin{pmatrix}1&1\\8&1\end{pmatrix}$, with $F\equiv A\pmod8$. The comparison target fails the same finite test.\par
This is local matrix algebra, not global automorphy.
}{A supplied two sided inverse certificate makes one finite congruence decisive in the stated ball.}

\GuideSpread{Modular singularity and genuine singularity differ}{The inverse denominator can exceed the scalar valuation}{22_singular.pdf}{
For the unipotent square center, the derivative is invertible over $\Q_2$, even though its reduction modulo 2 is singular. The inverse needs denominator $2^2$, so $s=2>\nu_2(2)=1$.\par
The strict condition for the previous contraction theorem is $m>s$. Replacing $s$ by the scalar degree valuation would enlarge the claimed ball without justification.\par
At $A=\operatorname{diag}(1,-1)$, the square derivative multiplies the diagonal directions by 2 and $-2$, and kills both off diagonal directions. It is singular over $\Q_2$.\par
That center cannot satisfy the same two sided inverse certificate.
}{Audit the full matrix derivative before applying an inverse based lifting theorem.}

\GuideSpread{A tower is a chain of compatible choices}{Finite roots illustrate why arbitrary selections can fail}{23_tower.pdf}{
The tree contains every odd root of $x^2=1$ modulo 2, 4, 8, and 16. Each edge is reduction to the preceding level. All node values and edges are exact.\par
The highlighted branch chooses 1 at every level. The root 3 modulo 8 has no successor modulo 16, although both levels have solutions. Thus a root chosen at one level need not extend.\par
For compact finite solution sets with compatible reduction maps, nonempty levels give an infinite compatible branch by the usual finite branching argument. That conclusion requires the maps and the solution conditions to agree across levels.\par
A local matrix root already constructed in $\Z_p$ automatically supplies its compatible reductions.
}{Record transition maps and compatible choices, not just separate existence flags.}

\GuideSpread{The remaining attachment map is concrete}{Generic algebraic tools and arithmetic existence have different jobs}{24_open_gate.pdf}{
The upper route begins with an actual finite Satake automorphic descent. A residual carrier at a good separating prime, one semisimple l adic realization, or a suitable quotient prime lifting theorem can close the cyclic attachment step under its hypotheses.\par
The lower route records the generic local toolkit: an actual deformation problem, the correct cochain comparison, compatible defects, correction exactness, and controlled lifting. The square zero and Hensel results apply once those inputs are proved.\par
The dashed connectors are the arithmetic application theorems still needed. A formal datum equating two spaces does not construct that comparison. An auxiliary root does not construct an automorphic carrier.\par
Neither route is depicted as completed Strong Artin.
}{The open theorem must construct the arithmetic input and verify the attachment it claims.}
